\chapter{Differential Equations and Dynamical Systems}

\section{First-Order Differential Equations and Initial Value Problems}

\subsection{Differential Equations and Families of Solutions}

An ordinary differential equation of first order is an equation that involves an unknown function of a single variable, which we shall generally denote by $y = y(t)$, together with its first derivative. The general form of such an equation is
\[
y' = F(t,y),
\]
where $F(t,y)$ is a given expression in the two variables $t$ and $y$. Several important classes of first-order equations arise as special cases of this general form, according to the specific dependence of $F$ on $t$ and $y$. A pure-time equation is one of the form $y' = f(t)$, in which the right-hand side depends only on the independent variable. A separable equation is one of the form $y' = a(t) g(y)$, in which the right-hand side factors as a product of a function of $t$ alone and a function of $y$ alone. An autonomous equation is the special case of a separable equation in which the dependence on $t$ disappears entirely, so that $y' = g(y)$. Finally, a linear equation is one of the form $y' = a(t) y + f(t)$, in which the right-hand side is affine in the unknown $y$.

\begin{definition}
A function $y = y(t)$ is called a solution of the differential equation $y' = F(t,y)$ on the interval $I$ if the equation is satisfied identically, that is, if
\[
y'(t) = F(t, y(t))
\]
for every $t \in I$.
\end{definition}

As a simple illustration of this definition, one may verify directly that $y(t) = 10 e^{2t}$ is a solution of the equation $y' = 2y$, since differentiating this expression indeed yields
\[
\begin{aligned}
y'(t) &= 20 e^{2t} \\
&= 2 y(t).
\end{aligned}
\]

When asked to solve a differential equation, one is generally expected to determine all possible solutions of the equation, rather than a single particular solution; this immediately raises the natural question of how many solutions a first-order equation may possess.


To address this question, it is instructive to consider the simplest possible class of first-order equations, namely the pure-time equations $y' = f(t)$. For an arbitrary constant $C \in \mathbb{R}$, the solutions are given by
\[
\begin{aligned}
y(t) &= \int f(t) \, dt \\
&= g(t) + C,
\end{aligned}
\]
where $g$ denotes any fixed antiderivative of $f$; the equation therefore possesses infinitely many solutions, parametrized by the arbitrary constant $C$. This phenomenon is not peculiar to pure-time equations, but is in fact a general feature of first-order ordinary differential equations: such equations typically possess infinitely many solutions, depending on a single free parameter, which may be fixed by imposing an additional condition, as will be discussed below in connection with initial value problems.

Figure~\ref{fig:atlas-ode-direction-field} relates a differential equation to its direction field and family of solutions.

\begin{figure}[!htbp]
\centering
\begin{tikzpicture}[>=Stealth]
\begin{axis}[
  width=12.5cm,height=6.5cm,
  axis lines=left,
  ticklabel style={font=\small},label style={font=\small},
  xlabel={$t$},ylabel={$y$},
  xmin=0,xmax=2.1,ymin=0,ymax=3.3,
  xtick={0,0.5,1,1.5,2},ytick={0,1,2,3},
  legend style={at={(0.5,1.02)},anchor=south,legend columns=3,
    font=\scriptsize,draw=black!20,fill=white,column sep=1em}
]
\draw[black!25] (axis cs:-0.07500,0.00000) -- (axis cs:0.07500,0.00000);
\draw[black!25] (axis cs:-0.06708,0.46646) -- (axis cs:0.06708,0.53354);
\draw[black!25] (axis cs:-0.05303,0.94697) -- (axis cs:0.05303,1.05303);
\draw[black!25] (axis cs:-0.04160,1.43760) -- (axis cs:0.04160,1.56240);
\draw[black!25] (axis cs:-0.03354,1.93292) -- (axis cs:0.03354,2.06708);
\draw[black!25] (axis cs:-0.02785,2.43036) -- (axis cs:0.02785,2.56964);
\draw[black!25] (axis cs:-0.02372,2.92885) -- (axis cs:0.02372,3.07115);
\draw[black!25] (axis cs:0.17500,0.00000) -- (axis cs:0.32500,0.00000);
\draw[black!25] (axis cs:0.18292,0.46646) -- (axis cs:0.31708,0.53354);
\draw[black!25] (axis cs:0.19697,0.94697) -- (axis cs:0.30303,1.05303);
\draw[black!25] (axis cs:0.20840,1.43760) -- (axis cs:0.29160,1.56240);
\draw[black!25] (axis cs:0.21646,1.93292) -- (axis cs:0.28354,2.06708);
\draw[black!25] (axis cs:0.22215,2.43036) -- (axis cs:0.27785,2.56964);
\draw[black!25] (axis cs:0.22628,2.92885) -- (axis cs:0.27372,3.07115);
\draw[black!25] (axis cs:0.42500,0.00000) -- (axis cs:0.57500,0.00000);
\draw[black!25] (axis cs:0.43292,0.46646) -- (axis cs:0.56708,0.53354);
\draw[black!25] (axis cs:0.44697,0.94697) -- (axis cs:0.55303,1.05303);
\draw[black!25] (axis cs:0.45840,1.43760) -- (axis cs:0.54160,1.56240);
\draw[black!25] (axis cs:0.46646,1.93292) -- (axis cs:0.53354,2.06708);
\draw[black!25] (axis cs:0.47215,2.43036) -- (axis cs:0.52785,2.56964);
\draw[black!25] (axis cs:0.47628,2.92885) -- (axis cs:0.52372,3.07115);
\draw[black!25] (axis cs:0.67500,0.00000) -- (axis cs:0.82500,0.00000);
\draw[black!25] (axis cs:0.68292,0.46646) -- (axis cs:0.81708,0.53354);
\draw[black!25] (axis cs:0.69697,0.94697) -- (axis cs:0.80303,1.05303);
\draw[black!25] (axis cs:0.70840,1.43760) -- (axis cs:0.79160,1.56240);
\draw[black!25] (axis cs:0.71646,1.93292) -- (axis cs:0.78354,2.06708);
\draw[black!25] (axis cs:0.72215,2.43036) -- (axis cs:0.77785,2.56964);
\draw[black!25] (axis cs:0.72628,2.92885) -- (axis cs:0.77372,3.07115);
\draw[black!25] (axis cs:0.92500,0.00000) -- (axis cs:1.07500,0.00000);
\draw[black!25] (axis cs:0.93292,0.46646) -- (axis cs:1.06708,0.53354);
\draw[black!25] (axis cs:0.94697,0.94697) -- (axis cs:1.05303,1.05303);
\draw[black!25] (axis cs:0.95840,1.43760) -- (axis cs:1.04160,1.56240);
\draw[black!25] (axis cs:0.96646,1.93292) -- (axis cs:1.03354,2.06708);
\draw[black!25] (axis cs:0.97215,2.43036) -- (axis cs:1.02785,2.56964);
\draw[black!25] (axis cs:0.97628,2.92885) -- (axis cs:1.02372,3.07115);
\draw[black!25] (axis cs:1.17500,0.00000) -- (axis cs:1.32500,0.00000);
\draw[black!25] (axis cs:1.18292,0.46646) -- (axis cs:1.31708,0.53354);
\draw[black!25] (axis cs:1.19697,0.94697) -- (axis cs:1.30303,1.05303);
\draw[black!25] (axis cs:1.20840,1.43760) -- (axis cs:1.29160,1.56240);
\draw[black!25] (axis cs:1.21646,1.93292) -- (axis cs:1.28354,2.06708);
\draw[black!25] (axis cs:1.22215,2.43036) -- (axis cs:1.27785,2.56964);
\draw[black!25] (axis cs:1.22628,2.92885) -- (axis cs:1.27372,3.07115);
\draw[black!25] (axis cs:1.42500,0.00000) -- (axis cs:1.57500,0.00000);
\draw[black!25] (axis cs:1.43292,0.46646) -- (axis cs:1.56708,0.53354);
\draw[black!25] (axis cs:1.44697,0.94697) -- (axis cs:1.55303,1.05303);
\draw[black!25] (axis cs:1.45840,1.43760) -- (axis cs:1.54160,1.56240);
\draw[black!25] (axis cs:1.46646,1.93292) -- (axis cs:1.53354,2.06708);
\draw[black!25] (axis cs:1.47215,2.43036) -- (axis cs:1.52785,2.56964);
\draw[black!25] (axis cs:1.47628,2.92885) -- (axis cs:1.52372,3.07115);
\draw[black!25] (axis cs:1.67500,0.00000) -- (axis cs:1.82500,0.00000);
\draw[black!25] (axis cs:1.68292,0.46646) -- (axis cs:1.81708,0.53354);
\draw[black!25] (axis cs:1.69697,0.94697) -- (axis cs:1.80303,1.05303);
\draw[black!25] (axis cs:1.70840,1.43760) -- (axis cs:1.79160,1.56240);
\draw[black!25] (axis cs:1.71646,1.93292) -- (axis cs:1.78354,2.06708);
\draw[black!25] (axis cs:1.72215,2.43036) -- (axis cs:1.77785,2.56964);
\draw[black!25] (axis cs:1.72628,2.92885) -- (axis cs:1.77372,3.07115);
\draw[black!25] (axis cs:1.92500,0.00000) -- (axis cs:2.07500,0.00000);
\draw[black!25] (axis cs:1.93292,0.46646) -- (axis cs:2.06708,0.53354);
\draw[black!25] (axis cs:1.94697,0.94697) -- (axis cs:2.05303,1.05303);
\draw[black!25] (axis cs:1.95840,1.43760) -- (axis cs:2.04160,1.56240);
\draw[black!25] (axis cs:1.96646,1.93292) -- (axis cs:2.03354,2.06708);
\draw[black!25] (axis cs:1.97215,2.43036) -- (axis cs:2.02785,2.56964);
\draw[black!25] (axis cs:1.97628,2.92885) -- (axis cs:2.02372,3.07115);
\addplot[figblue,very thick,domain=0:2,samples=100] {0.2*exp(x)};
\addlegendentry{$C=0.2$}
\addplot[figred,thick,domain=0:2,samples=100] {0.3*exp(x)};
\addlegendentry{$C=0.3$}
\addplot[figgreen,thick,domain=0:2,samples=100] {0.4*exp(x)};
\addlegendentry{$C=0.4$}
\end{axis}
\end{tikzpicture}
\caption{Short segments indicate the slope $y$ prescribed by $y^{\prime}=y$. The curves are exact solutions $y(t)=Ce^t$ for three positive initial values. Each initial value selects one curve; the direction field itself represents all such solutions, not a single trajectory.}
\label{fig:atlas-ode-direction-field}
\end{figure}

\subsection{Techniques for Solving First-Order Equations}

Among the general theory of first-order equations, explicit techniques are available for finding a formula for the general solution, also called the general integral, of several important classes: pure-time equations $y' = f(t)$, separable equations $y' = a(t) g(y)$, including the autonomous case, and linear equations
\[
y' = a(t) y + f(t).
\]

We now illustrate the method applicable to separable equations through two worked examples, before stating the general procedure.

\begin{example}
Consider the equation $y' = \varepsilon y$, where $\varepsilon$ is a given constant. One immediately identifies the constant solution $y(t) \equiv 0$. Setting $y \neq 0$ and dividing by $y$, the equation becomes
\[
\frac{y'(t)}{y(t)} = \varepsilon \implies \int \frac{y'(t)}{y(t)} \, dt = \int \varepsilon \, dt.
\]

The integral on the left-hand side is elementary, yielding $\ln |y(t)| = \varepsilon t + c$ for an arbitrary constant $c$; setting $C = e^c$, we obtain $|y(t)| = C e^{\varepsilon t}$. This furnishes a branch of positive solutions $y(t) = C e^{\varepsilon t}$, for $C > 0$, together with a branch of negative solutions $y(t) = -C e^{\varepsilon t}$, again for $C > 0$; together with the constant solution $y \equiv 0$ found initially, these exhaust all solutions of the equation.
\end{example}

\begin{example}
Consider the equation $y' = t y^3$. As before, one immediately identifies the constant solution $y(t) \equiv 0$. Setting $y \neq 0$ and separating the variables by dividing by $y^3$, the equation becomes
\[
\frac{y'(t)}{y^3(t)} = t \implies \int \frac{y'(t)}{y^3(t)} \, dt = \int t \, dt.
\]
The integral on the left-hand side is again elementary, yielding
\[
-(1/(2 y^2(t))) = (t^2/2) + c;
\]
setting $C = -2c$, we obtain two branches of solutions,
\[
y(t) = \frac{1}{\sqrt{C - t^2}}, \qquad y(t) = -\frac{1}{\sqrt{C - t^2}},
\]
each defined on the maximal interval $(-\sqrt{C}, \sqrt{C})$, for every fixed $C > 0$.
\end{example}

These two examples illustrate the general method applicable to the broader class of separable equations, defined as follows.

\begin{definition}[Separable Equation]
An equation of the form $y' = a(t) g(y)$ is called a separable equation. The special case in which $a(t)$ is constant, so that the equation takes the form $y' = g(y)$, is called an autonomous equation.
\end{definition}

Assuming that $a$ and $g$ are continuous functions, the general solution of a separable equation $y' = a(t) g(y)$ can be obtained by means of a systematic four-step procedure. The first step consists of searching for equilibrium solutions, if any exist, namely constant solutions of the form $y(t) \equiv r \in \mathbb{R}$, which correspond precisely to the roots of the equation $g(r) = 0$. The second step consists of separating the two variables $t$ and $y$: assuming that $g(y) \neq 0$, one divides the equation by $g(y)$ and integrates both sides with respect to $t$, obtaining
\[
\int \frac{y'(t)}{g(y(t))} \, dt = \int a(t) \, dt.
\]

The third step consists of evaluating both integrals explicitly: denoting by $A(t)$ a given antiderivative of $a(t)$, and by $H(u)$ an antiderivative of $1/g(u)$, one obtains the implicit relation
\[
H(y(t)) = A(t) + C,
\]
for an arbitrary constant $C \in \mathbb{R}$. The fourth and final step consists of inverting this implicit relation, whenever $H$ is invertible, to obtain the explicit formula
\[
y(t) = H^{-1}(A(t) + C);
\]
it should be noted that restrictions on the domain of $t$, possibly depending on the value of $C$, may arise at this final step, as illustrated in the two examples above.

\subsection{Applications to Population Dynamics}


One of the most important applications of the theory of first-order differential equations lies in biology, where a wide variety of processes can be modeled by such equations, providing considerable insight into the dynamics of living organisms, that is, into how individuals and populations evolve over time. The earliest and simplest such model, proposed by Malthus in 1830, describes the growth of a population under idealized conditions.

Let $N(t)$ denote the number of individuals in a population at time $t$. The Malthus model rests on three assumptions: each individual produces offspring at a constant rate $\lambda$, so that the total birth rate at time $t$ equals $\lambda N(t)$; and individuals die at a constant per-capita rate $\mu$, so that the total loss rate through death at time $t$ equals $\mu N(t)$. Since the rate of change of the population size equals the total birth rate minus the total death rate, with $\varepsilon = \lambda - \mu$, one obtains the following equivalent forms of the differential equation:
\[
\begin{aligned}
N'(t) &= \lambda N(t) - \mu N(t), \\
\frac{dN}{dt} &= \varepsilon N.
\end{aligned}
\]

The constant $\varepsilon$ is called the biological potential, or Malthus parameter, of the population. This linear autonomous equation can be solved explicitly by the separation of variables technique developed above, yielding exponential growth or decay depending on the sign of $\varepsilon$. A comparison with experimental data on the growth of a population of yeast cells, for instance, shows reasonably good agreement between the model $dN/dt = 0.5\, N$, with initial condition $N(0) = 0.2 \times 10^6$ cells per milliliter, and the observed population counts, at least during the early stages of growth.


The Malthus model, while adequate for describing the early stages of population growth, fails to capture an important biological phenomenon known as the logistic effect, namely the intraspecific competition for limited resources: when the population density is high, one typically observes a decrease in the birth rate together with an increase in the mortality rate, effects that the constant biological potential of the Malthus model cannot represent. The logistic, or Verhulst, model, proposed in 1838, addresses this limitation by allowing the biological potential to depend on the population size itself.

Specifically, instead of assuming a constant biological potential, the logistic model assumes that the biological potential decreases linearly as the population size increases, starting from a value $\varepsilon > 0$ when $N = 0$ and reaching the value $0$ when $N = k$, where the constant $k$ is called the carrying capacity of the environment. This linear dependence is expressed by the formula
\[
\varepsilon(N) = \varepsilon \left( 1 - \frac{N}{k} \right),
\]
and substituting this expression for the biological potential into the growth equation yields the logistic equation
\[
\frac{dN}{dt} = \varepsilon \left( 1 - \frac{N}{k} \right) N.
\]

As with the Malthus model, a comparison with experimental data on the growth of a yeast population shows that the logistic model
\[
dN/dt = (0.55 - 0.0026\, N)\, N,
\]
again with initial condition $N(0) = 0.2 \times 10^6$ cells per milliliter, reproduces the observed data considerably more accurately than the simple Malthus model, particularly as the population approaches its carrying capacity. As exercises illustrating the qualitative behavior of the logistic equation, the reader may solve the equation $y' = y(2-y)$ explicitly by the method of separation of variables, and may likewise solve the closely related equation $y' = y(y-2)$, verifying in the process that some solutions of this second equation remain defined for all time, while others blow up in finite time. Figure~\ref{fig:ch7-logistic} connects the logistic trajectory with its phase-line description.

\begin{figure}[!htbp]
\centering
\begin{tikzpicture}
\begin{axis}[
  width=6.2cm,height=6.2cm, axis lines=left,
  xlabel={$t$}, ylabel={$S(t)$},
  xmin=0,xmax=10, ymin=0,ymax=5.6,
  ytick={5},yticklabels={$K$}, xtick=\empty,
  ]
  \addplot[figblue,very thick,domain=0:10,samples=150]{5/(1+9*exp(-0.9*x))};
  \draw[figred,dashed] (axis cs:0,5) -- (axis cs:10,5);
\end{axis}
\begin{axis}[
  at={(6.9cm,0)},
  width=6.2cm,height=6.2cm, axis lines=middle,
  xlabel={$S$}, ylabel={$\dot S$},
  xmin=0,xmax=5.6, ymin=-1,ymax=1.6,
  xtick={5},xticklabels={$K$}, ytick=\empty,
  ]
  \addplot[figgreen,very thick,domain=0:5.6,samples=100]{0.5*x*(1-x/5)};
  \node[circle,fill=white,draw=black,inner sep=1.3pt] at (axis cs:0,0) {};
  \node[below right,fill=white,inner sep=1pt,font=\scriptsize]
    at (axis cs:0.08,-0.12) {unstable};
  \node[circle,fill=black,inner sep=1.3pt] at (axis cs:5,0) {};
  \node[below left,fill=white,inner sep=1pt,font=\scriptsize]
    at (axis cs:4.95,-0.12) {stable};
  \draw[vecA,shorten >=1pt] (axis cs:1.0,0) -- (axis cs:2.2,0);
  \draw[vecA,shorten >=1pt] (axis cs:3.0,0) -- (axis cs:4.2,0);
  \draw[vecA,shorten >=1pt] (axis cs:5.55,0) -- (axis cs:5.15,0);
\end{axis}
\end{tikzpicture}
\caption[Logistic model: S-curve and phase line]{Left: solution $S(t)$ approaching
  the carrying capacity $K$ asymptotically. Right: phase portrait $\dot S$ vs $S$, with
  blue arrows indicating the flow direction; $S=0$ is unstable, $S=K$ is stable.}
\label{fig:ch7-logistic}
\end{figure}

\subsection{The Initial Value Problem}

In many applications arising from physical or biological problems, one is typically not interested in the entire family of solutions of a differential equation, but rather in the single, particular solution that satisfies an additional condition of the form $y(t_0) = y_0$ for prescribed values $t_0$ and $y_0$. Such a condition is called an initial condition, and the problem of finding the solution of the differential equation that satisfies a given initial condition is called an initial value problem, or Cauchy problem.

\begin{definition}[Cauchy Problem]
A Cauchy problem is a system consisting of a differential equation together with an initial condition, of the form
\[
\begin{cases}
y' = F(t,y), \\
y(t_0) = y_0,
\end{cases}
\]
where $t_0$ and $y_0$ are assigned values.
\end{definition}

A central theoretical question concerns whether a given Cauchy problem admits a solution, and if so, whether that solution is unique; a problem for which both properties hold is said to be well-posed.

\begin{definition}[Well-Posedness]
The Cauchy problem with initial datum $(t_0, y_0)$ is said to be well-posed if it admits a unique solution $y(\cdot)$ defined in a neighborhood of $t_0$.
\end{definition}

Well-posedness has an important geometric consequence: it implies that exactly one solution curve of the differential equation passes through the point $(t_0, y_0)$ in the $(t,y)$-plane, and consequently that two distinct solution curves can never intersect at a common point. In particular, if the Cauchy problem is well-posed for every admissible choice of initial condition, then a constant solution $y \equiv k$ can never be crossed by any other solution curve; the horizontal line $y = k$ therefore constitutes an impassable barrier for every other solution curve of the equation, a fact of considerable qualitative importance when sketching the family of solutions of a given equation.

For separable equations, well-posedness of the associated Cauchy problem can be guaranteed under mild regularity assumptions on the functions $a$ and $g$, though only in a local sense, as the following theorem makes precise.

\begin{theorem}[Local Well-Posedness for Separable Equations]
Let $a(t)$ be continuous on an interval $I = (\alpha, \beta)$, and let $g(y)$ be differentiable with $g'(y)$ continuous on an interval $J = (\gamma, \delta)$. Then, for every choice of $t_0 \in (\alpha,\beta)$ and every choice of initial datum $y_0 \in (\gamma,\delta)$, the Cauchy problem
\[
\begin{cases}
y' = a(t) g(y), \\
y(t_0) = y_0,
\end{cases}
\]
has a unique solution $y(t)$, defined on an interval $(t_0 - \delta_0, t_0 + \delta_0)$ for some $\delta_0 > 0$.
\end{theorem}

\begin{remark}
It is important to emphasize that the size of the interval of existence $\delta_0$ guaranteed by this theorem cannot, in general, be predicted theoretically from the hypotheses alone; the theorem asserts only the local existence and uniqueness of the solution, without providing quantitative information on how far this solution can be extended.
\end{remark}

The local nature of the well-posedness theorem stated above is not a mere technical artifact: it reflects genuine phenomena that can occur even for very simple separable equations, as the following two examples illustrate.

\begin{remark}
Existence of solutions is, in general, not global in time: there exist solutions of separable equations that are not defined on the whole interval $I$ where the coefficient $a(t)$ is continuous.
\end{remark}

\begin{example}[Blow-Up in Finite Time]
Consider the Cauchy problem $y' = y^2$, $y(0) = 1$. This equation is separable, with $a(t) \equiv 1$, continuous on all of $\mathbb{R}$, and $g(y) = y^2$. The problem admits a unique solution, given explicitly by $y(t) = 1/(1-t)$. Since
\[
\lim_{t \to 1^-} y(t) = +\infty,
\]
the solution blows up at the finite time $t=1$, and its maximal interval of existence is $I_{\max} = (-\infty, 1)$, strictly smaller than the interval $\mathbb{R}$ on which the coefficients of the equation are defined.
\end{example}

\begin{remark}
A lack of uniqueness may likewise occur whenever the hypotheses of the Local Well-Posedness Theorem, particularly the differentiability of $g$, are violated.
\end{remark}

\begin{example}[Lack of Uniqueness]
Consider the Cauchy problem
\[
y' = 3\sqrt[3]{y^2},
\]
$y(0) = 0$. This equation is separable, with
\[
g(y) = 3\sqrt[3]{y^2},
\]
which fails to be differentiable at $y = 0$. The null constant function $y(t) = 0$, for every $t \in \mathbb{R}$, is one solution of this problem; however, the function $y(t) = t^3$ provides another solution, defined globally on $\mathbb{R}$, distinct from the first.
\end{example}

In fact, the Cauchy problem of the preceding example admits infinitely many solutions, a phenomenon known as Peano's brush (\emph{Pennello di Peano}); as a further illustration, the function
\[
y(t) =
\begin{cases}
0 & t \leq 3, \\
(t-3)^3 & t > 3,
\end{cases}
\]
provides yet another solution of the same Cauchy problem, distinct from both of those exhibited above.

As further exercises illustrating the qualitative theory of separable equations, the reader may consider the equation $y' = t y^3$ and determine, for each real number $a$, the solution of the initial value problem with condition $y(1) = a$; observing that the size of the maximal interval of existence depends sensitively on the specific initial datum chosen, one may plot several of these solutions on a common graph in order to visualize this dependence. A closely related exercise concerns the equation
\[
y' = (1/2)(y^2 - 1),
\]
for which the solutions corresponding to the initial condition $y(0) = a$, again depending on the real parameter $a$, may be found explicitly and compared.

Figure~\ref{fig:atlas-finite-time-blowup} shows the finite right endpoint of a maximal solution interval.

\begin{figure}[!htbp]
\centering
\begin{tikzpicture}[>=Stealth]
\begin{axis}[
  width=12.3cm,height=6.3cm,
  axis lines=left,
  ticklabel style={font=\small},label style={font=\small},
  xlabel={$t$},ylabel={$y(t)$},
  xmin=-1,xmax=1.25,ymin=0,ymax=10.5,
  xtick={-1,0,1},ytick={1,3,5,7,9}
]
\addplot[figblue,very thick,-{Stealth[length=2.4mm]},domain=-1:0.89,samples=180] {1/(1-x)};
\draw[figred,dashed,thick] (axis cs:1,0) -- (axis cs:1,9.8);
\node[figred,above left,fill=white,inner sep=1pt,font=\small]
  at (axis cs:0.98,9.6) {$T_+=1$};
\addplot[only marks,figblue,mark=*] coordinates {(0,1)};
\node[figblue,above left,font=\scriptsize] at (axis cs:0,1) {$y(0)=1$};
\end{axis}
\end{tikzpicture}
\caption{The solution $y(t)=1/(1-t)$ with initial value $y(0)=1$ satisfies $y^{\prime}=y^2$ and is defined on $(-\infty,1)$. The vertical dashed line marks the blow-up time, not part of the solution. The equation has smooth coefficients, but this trajectory cannot be continued through that time with a finite value.}
\label{fig:atlas-finite-time-blowup}
\end{figure}

\section{Linear First-Order Equations}

\subsection{Global Well-Posedness}

We now turn to the class of linear first-order equations, already introduced above, which take the general form $y' = a(t) y + f(t)$. Several familiar equations from physics and other applications fall within this class, including the equation $P' = 2P$ governing simple population or capital growth, the equation $y' = y + 3x$, and the equation $y' - (1/t) y = \ln t$.

Unlike the case of general separable equations, for which well-posedness of the Cauchy problem could only be guaranteed locally in time, linear equations enjoy a considerably stronger property: well-posedness holds globally on the entire interval where the coefficients are defined. This distinctive feature, a direct consequence of the linear structure of the equation, is recorded in the following theorem.

\begin{theorem}[Global Well-Posedness for Linear Equations]
Let $I$ be an interval, and assume that $a(t)$ and $f(t)$ are both continuous on $I$. Then, for every choice of initial time $t_0 \in I$ and every choice of initial value $y_0 \in \mathbb{R}$, the Cauchy problem
\[
\begin{cases}
y' = a(t) y + f(t), \\
y(t_0) = y_0,
\end{cases}
\]
admits a unique solution $y(t)$, defined on the whole interval $I$.
\end{theorem}


\subsection{The Integrating Factor Method}

Every first-order linear differential equation can be solved explicitly by means of a systematic technique known as the integrating factor method. As a first illustration, one may consider the equation $y' + (1/x) y = 2$, which can be solved by multiplying both sides by the factor $x$; this particular example already suggests the general strategy, which we now describe in full.

To solve the linear differential equation written in the form $y' - a(t) y = f(t)$, one multiplies both sides by the integrating factor
\[
e^{-\int a(t) \, dt}
\]
and subsequently integrates both sides with respect to $t$.

\begin{proof}
Let $A(t)$ denote any fixed antiderivative of $a(t)$. Multiplying the equation
\[
y' - a(t) y = f(t)
\]
by the integrating factor $e^{-A(t)}$, we obtain
\begin{equation*}
e^{-A(t)} \big[ y' - a(t) y \big] = e^{-A(t)} f(t). \tag{$*$}
\end{equation*}

Consider now the product function $y(t) e^{-A(t)}$; its derivative, computed by the product rule together with the Chain Rule applied to $e^{-A(t)}$, is precisely
\[
\frac{d}{dt} \big[ y(t) e^{-A(t)} \big] = e^{-A(t)} \big[ y'(t) - a(t) y(t) \big],
\]
so that equation $(*)$ may be rewritten as
\[
\frac{d}{dt} \big[ y(t) e^{-A(t)} \big] = e^{-A(t)} f(t).
\]

By the definition of the indefinite integral, this identity gives the following relation with an arbitrary constant $c \in \mathbb{R}$:
\[
y(t) e^{-A(t)} + c = \int e^{-A(t)} f(t) \, dt.
\]
Equivalently, the solution is
\[
y(t) = C \, e^{A(t)} + e^{A(t)} \int e^{-A(t)} f(t) \, dt,
\]
where $C$ denotes an arbitrary real constant, absorbing the constant $-c$ from the previous line. This establishes the desired formula for the general solution.
\end{proof}

The proof just given yields an explicit formula for the general solution of any linear first-order equation with continuous coefficients, which we now record for reference.

\begin{theorem}[General Integral of the Linear Equation]
Let $a(t)$ and $f(t)$ be continuous functions on an interval $I$, and let $A(t)$ denote an antiderivative of $a$ on $I$. For an arbitrary constant $C \in \mathbb{R}$, the solutions of $y' = a(t) y + f(t)$ on $I$ and a particular solution are given by
\[
\begin{aligned}
y(t) &= C \, e^{A(t)} + y_P(t), \\
y_P(t) &= e^{A(t)} \int e^{-A(t)} f(t) \, dt.
\end{aligned}
\]

\end{theorem}

\begin{remark}
The general solution formula obtained above can be written suggestively as
\[
y(t) = y_{\text{homo}}(t) + y_P(t).
\]
The homogeneous term is
\[
y_{\text{homo}}(t) = C e^{A(t)}.
\]

This is the general solution of the associated homogeneous equation $y' = a(t) y$, obtained by setting $f \equiv 0$, and $y_P(t)$ denotes a single particular solution of the full nonhomogeneous equation. This decomposition anticipates a structural principle that will play an even more prominent role in the theory of second-order linear equations developed in the following section.
\end{remark}

\section{Second-Order Linear Differential Equations}

\subsection{Second-Order Linear Equations and Initial Data}

We now turn to the study of second-order linear differential equations, which take the general form
\[
a(t) y'' + b(t) y' + c(t) y = f(t), \qquad t \in I,
\]
where $a$, $b$, and $c$ are continuous functions on the interval $I$, called the coefficient functions, subject to the nondegeneracy condition $a(t) \neq 0$ for every $t \in I$, and where $f$ is a continuous function on $I$, called the forcing term. Equations of this type arise naturally in a wide variety of physical contexts; a particularly simple example is furnished by the equation of motion of a falling body, $m y'' = -mg$, whose general solution is readily found by direct integration to be
\[
y(t) = -(1/2) g t^2 + v_0 t + y_0,
\]
where $v_0$ and $y_0$ denote the initial velocity and initial position of the body, respectively. As for first-order linear equations, the associated initial value problem for a second-order linear equation is well-posed globally on the interval of continuity of the coefficients, subject now to the specification of two initial conditions, since the equation involves a second derivative.

\begin{theorem}[Cauchy Theorem for Second-Order Linear Equations]
Assume that $a$, $b$, $c$, and $f$ are continuous functions on an interval $I$, with $a(t) \neq 0$ throughout $I$. Then, for any given $t_0 \in I$, $y_0 \in \mathbb{R}$, and $v_0 \in \mathbb{R}$, the Cauchy problem
\[
\begin{cases}
a(t) y'' + b(t) y' + c(t) y = f(t), \\
y(t_0) = y_0, \\
y'(t_0) = v_0,
\end{cases}
\]
has a unique solution $y : I \to \mathbb{R}$.
\end{theorem}

\subsection{The Homogeneous Equation}

We first study the homogeneous equation associated with the general second-order linear equation, namely
\begin{equation*}
a(t) y'' + b(t) y' + c(t) y = 0, \qquad t \in I, \tag{$\star$}
\end{equation*}
where, as above, $a$, $b$, $c$ are continuous on $I$ with $a(t) \neq 0$ throughout. The linear structure of this equation is reflected in the following superposition principle, which shows that the set of solutions is closed under linear combination.

\begin{theorem}[Superposition Principle for the Homogeneous Equation]
If $y_1(t)$ and $y_2(t)$ are both solutions of the linear homogeneous equation $(\star)$ on the interval $I$, and $c_1, c_2$ are arbitrary constants, then the function
\[
y(t) = c_1 y_1(t) + c_2 y_2(t)
\]
is also a solution of $(\star)$ on $I$.
\end{theorem}

This superposition principle, together with a suitable notion of linear independence, allows one to describe the entire family of solutions of the homogeneous equation in terms of just two particular solutions.

\begin{theorem}[General Integral of the Homogeneous Equation]
If $y_1(t)$ and $y_2(t)$ are linearly independent solutions of the linear homogeneous equation $(\star)$ on the interval $I$, meaning that the ratio $y_1/y_2$ is not constant on $I$, then the general solution of $(\star)$ is given by
\[
y(t) = c_1 y_1(t) + c_2 y_2(t), \qquad t \in I,
\]
where $c_1$ and $c_2$ are arbitrary constants.
\end{theorem}

This result shows that, once two particular linearly independent solutions of the homogeneous equation are known, every solution of the equation is thereby determined; equivalently, one says that the set of all solutions of the second-order linear homogeneous equation $(\star)$ forms a vector space of dimension two.

In general, discovering two particular solutions of a second-order linear equation with variable coefficients can be a genuinely difficult problem. It is, however, always possible to do so explicitly when the coefficients are constant, that is, for an equation of the form $a y'' + b y' + c y = 0$ with $a, b, c \in \mathbb{R}$ and $a \neq 0$; the question of how to find two linearly independent solutions of such a constant-coefficient homogeneous equation, according to the sign of the discriminant of the associated characteristic equation, is addressed through dedicated techniques developed elsewhere in the course.

\subsection{The Nonhomogeneous Equation}

We now turn to the general, nonhomogeneous second-order linear equation,
\begin{equation*}
a(t) y'' + b(t) y' + c(t) y = f(t), \qquad t \in I, \tag{$\star\star$}
\end{equation*}
where $a$, $b$, $c$ are continuous on $I$, with $a(t) \neq 0$ throughout, and $f$ is a given continuous external forcing term. As in the homogeneous case, the linear structure of the equation gives rise to a corresponding superposition principle, which in this setting relates solutions of the equation with different forcing terms.

\begin{theorem}[Superposition Principle for the Nonhomogeneous Equation]
If $y_1(t)$ and $y_2(t)$ are solutions of the linear equation $(\star\star)$ with external forcing terms $f_1$ and $f_2$ respectively, and $c_1, c_2$ are arbitrary constants, then the function
\[
y(t) = c_1 y_1(t) + c_2 y_2(t)
\]
is a solution of $(\star\star)$ with forcing term
\[
f(t) = c_1 f_1(t) + c_2 f_2(t).
\]

\end{theorem}

The general solution of the nonhomogeneous equation can be described in terms of the solutions of the associated homogeneous equation, together with any single particular solution of the nonhomogeneous equation itself, according to the following theorem.

\begin{theorem}[General Integral of the Nonhomogeneous Equation]
The general solution of the nonhomogeneous differential equation $(\star\star)$ can be written as $y = y_{\text{homo}} + y_P$, where $y_P$ is any particular solution of $(\star\star)$, and $y_{\text{homo}}$ denotes the general solution of the associated complementary homogeneous equation $(\star)$.
\end{theorem}

This theorem shows that, once two linearly independent solutions of the associated homogeneous equation $(\star)$ are known, together with a single particular solution of the nonhomogeneous equation $(\star\star)$, every solution of $(\star\star)$ is thereby determined, exactly as in the analogous situation encountered for first-order linear equations.

As in the homogeneous case, finding a particular solution of a nonhomogeneous second-order linear equation with variable coefficients can be genuinely difficult in general. When the coefficient functions are constant, however, that is, for an equation of the form
\[
a y'' + b y' + c y = f(t)
\]
with $a, b, c \in \mathbb{R}$ and $a \neq 0$, a systematic technique known as the Method of Undetermined Coefficients is available, which allows one to guess the form of a particular solution directly from the structure of the forcing term $f(t)$, whether polynomial, exponential, or trigonometric in nature; the detailed development of this method, together with the classification of the homogeneous solutions according to the sign of the discriminant $\Delta$ of the characteristic equation, is left to the dedicated material provided for further study.

\section{Existence, Uniqueness, and Maximal Solutions}

\subsection{Initial Data and Local Lipschitz Regularity}

We consider an autonomous system of ODEs in $\R^N$:
\begin{equation}
\dot{\bv} = f(\bv), \quad \bv(0) = \bv_0
\end{equation}
where $f: \R^N \to \R^N$ is a vector field. A \textit{solution} to this initial value problem is a differentiable function $\bv: I \to \R^N$, defined on some interval $I$ containing $0$, such that
\[
(\dd/\dd t)\bv(t) = f(\bv(t))
\]
for all $t \in I$, and $\bv(0) = \bv_0$. To guarantee that solutions exist and are unique, we must impose regularity conditions on $f$. The standard condition is local Lipschitz continuity.

\begin{definition}[Locally Lipschitz Continuous]
A function $f: \R^N \to \R^N$ is \textit{locally Lipschitz continuous} if for every point $\bx \in \R^N$, there exists a neighborhood $U$ of $\bx$ and a constant $L_U > 0$ such that for all $\bx, \by \in U$:
\begin{equation}
\|f(\bx) - f(\by)\| \leq L_U \|\bx - \by\|
\end{equation}
where $\|\cdot\|$ denotes the standard Euclidean norm.
\end{definition}

\begin{remark}[Why the Lipschitz Condition?]
Continuity of $f$ alone (via Peano's existence theorem) guarantees that a solution exists, but it does not guarantee that the solution is \textit{unique}. Without uniqueness, trajectories could branch, making deterministic prediction impossible. The Lipschitz condition bounds the rate at which $f$ can change, preventing the vector field from becoming too steep and forcing trajectories to merge or branch.
\end{remark}

\subsection{Gronwall's Integral Inequality}


To prove uniqueness over the entire interval $[0, T]$, we rely on a fundamental integral inequality.

\begin{lemma}[Gronwall's Inequality]
Let $u$ and $\beta$ be continuous functions with the following domain and codomain:
\[
u, \beta: [0, T] \to [0, \infty).
\]
Suppose there exists a constant $C \geq 0$ such that the following bound holds for every $t \in [0, T]$:
\begin{equation}
u(t) \leq C + \int_0^t \beta(s) u(s) \dd s.
\end{equation}
For every $t$ in the same interval, the resulting estimate is
\begin{equation}
u(t) \leq C \exp\left( \int_0^t \beta(s) \dd s\right).
\end{equation}

\end{lemma}

\begin{proof}
We construct an explicit upper bound using the method of integrating factors.
Define the function $R(t)$ as the right-hand side of the inequality:
\begin{equation}
R(t) = C + \int_0^t \beta(s) u(s) \dd s
\end{equation}

By the given hypothesis, we immediately have $u(t) \leq R(t)$ for all $t$.
Notice that $R(0) = C$. By the Fundamental Theorem of Calculus, since the integrand $\beta(s)u(s)$ is continuous, $R(t)$ is differentiable and its derivative is:
\begin{equation}
R'(t) = \beta(t) u(t)
\end{equation}
Since $u(t) \leq R(t)$ and $\beta(t) \geq 0$, we can substitute to obtain a differential inequality for $R(t)$:
\begin{equation}
R'(t) \leq \beta(t) R(t) \implies R'(t) - \beta(t) R(t) \leq 0
\end{equation}
To solve this differential inequality, we multiply both sides by the integrating factor
\[
e^{-\int_0^t \beta(s) \dd s},
\]
which is strictly positive. Using the product rule for differentiation, we recognize the left side as the exact derivative of a product:
\begin{equation}
\begin{aligned}
\frac{\dd}{\dd t} \left[ R(t) e^{-\int_0^t \beta(s) \dd s} \right] &= R'(t) e^{-\int_0^t \beta} + R(t) \left( -\beta(t) e^{-\int_0^t \beta} \right) \\
&= e^{-\int_0^t \beta} \Big( R'(t) - \beta(t) R(t) \Big)
\end{aligned}
\end{equation}
The signs of the two factors are
\[
\begin{gathered}
e^{-\int_0^t \beta} > 0 \\
R'(t) - \beta(t) R(t) \leq 0,
\end{gathered}
\]
Hence the derivative of the product is non-positive:
\begin{equation}
\frac{\dd}{\dd t} \left[ R(t) e^{-\int_0^t \beta(s) \dd s} \right] \leq 0
\end{equation}
This implies that the function
\[
g(t) = R(t) e^{-\int_0^t \beta(s) \dd s}
\]
is monotonically non-increasing. Therefore, for any $t \geq 0$, $g(t) \leq g(0)$. Evaluating at $t=0$:
\begin{equation}
\begin{aligned}
g(0) &= R(0) e^0 \\
&= C \cdot 1 \\
&= C
\end{aligned}
\end{equation}
Thus,
\[
R(t) e^{-\int_0^t \beta(s) \dd s} \leq C.
\]
Multiplying both sides by the positive quantity
\[
e^{\int_0^t \beta(s) \dd s}
\]
yields:
\begin{equation}
R(t) \leq C \exp\left( \int_0^t \beta(s) \dd s \right)
\end{equation}
Finally, since $u(t) \leq R(t)$, we conclude that
\[
u(t) \leq C \exp\left( \int_0^t \beta(s) \dd s \right).
\]

\end{proof}

\subsection{Application of Gronwall's Inequality to Global Uniqueness}

\begin{theorem}[Global Uniqueness]
If $\bv_1(t)$ and $\bv_2(t)$ are two solutions to $\dot{\bv} = f(\bv)$ on the interval $[0, T]$ satisfying the same initial condition
\[
\begin{aligned}
\bv_1(0) &= \bv_2(0) \\
&= \bv_0,
\end{aligned}
\]
and if $f$ is globally Lipschitz continuous with constant $L$ (i.e.,
\[
\|f(\bx) - f(\by)\| \leq L\|\bx - \by\|
\]
for all $\bx, \by \in \R^N$), then $\bv_1(t) = \bv_2(t)$ for all $t \in [0, T]$.
\end{theorem}

\begin{proof}
First, we convert the differential equation into an integral equation. Since $\bv_i(t)$ is a solution,
\[
(\dd/\dd t)\bv_i(t) = f(\bv_i(t)).
\]
Integrating both sides from $0$ to $t$ and applying the Fundamental Theorem of Calculus yields:
\begin{equation}
\bv_i(t) - \bv_i(0) = \int_0^t f(\bv_i(s)) \dd s \implies \bv_i(t) = \bv_0 + \int_0^t f(\bv_i(s)) \dd s
\end{equation}
Now, consider the difference between the two solutions. Let
\[
u(t) = \|\bv_1(t) - \bv_2(t)\|.
\]
Subtracting the integral equations:
\begin{equation}
\bv_1(t) - \bv_2(t) = \int_0^t \Big( f(\bv_1(s)) - f(\bv_2(s)) \Big) \dd s
\end{equation}
Taking the norm of both sides and applying the triangle inequality for integrals,
\[
\|\int \mathbf{g}\| \leq \int \|\mathbf{g}\|
\]
gives:
\begin{equation}
\begin{aligned}
u(t) &= \left\| \int_0^t \Big( f(\bv_1(s)) - f(\bv_2(s)) \Big) \dd s \right\| \\
&\leq \int_0^t \left\| f(\bv_1(s)) - f(\bv_2(s)) \right\| \dd s
\end{aligned}
\end{equation}
We now explicitly apply the Lipschitz condition
\[
\|f(\bx) - f(\by)\| \leq L\|\bx - \by\|
\]
to the integrand:
\begin{equation}
\begin{aligned}
u(t) &\leq \int_0^t L \|\bv_1(s) - \bv_2(s)\| \dd s \\
&= L \int_0^t u(s) \dd s
\end{aligned}
\end{equation}

This is exactly the hypothesis of Gronwall's Inequality with $C = 0$ and $\beta(s) = L$. Applying the lemma:
\begin{equation}
\begin{aligned}
u(t) &\leq 0 \cdot \exp\left( \int_0^t L \dd s \right) \\
&= 0
\end{aligned}
\end{equation}
By the definition of a norm, the difference satisfies
\[
\begin{aligned}
u(t) &= \|\bv_1(t) - \bv_2(t)\| \\
&\geq 0
\end{aligned}
\]
Consequently, we must have $u(t) = 0$ for all $t \in [0, T]$. Therefore, $\bv_1(t) = \bv_2(t)$.
\end{proof}

\subsection{Maximal Interval of Existence}

Solutions to ODEs are guaranteed to exist locally. By continuously patching together local solutions, we can extend the solution until it either exists for all time or ``blows up''.

\begin{definition}[Maximal Interval of Existence]
The \textit{maximal interval of existence} $(T_-, T_+)$ is the largest open interval on which a solution $\bv(t)$ to the initial value problem is defined. A solution is called \textit{global} if $T_+ = \infty$ and $T_- = -\infty$.
\end{definition}

\begin{theorem}[Escape Lemma (Finite-Time Blow-up)]
Let $\bv(t)$ be the maximal solution defined on $(T_-, T_+)$. If the forward maximal time is finite, i.e., $T_+ < \infty$, then the trajectory must escape to infinity:
\begin{equation}
\lim_{t \to T_+^-} \|\bv(t)\| = \infty
\end{equation}

\end{theorem}

\begin{proof}
We prove this by contradiction. Suppose $T_+ < \infty$ but the trajectory does \textit{not} escape to infinity. This means there exists a sequence of times $t_n \to T_+^-$ such that $\|\bv(t_n)\|$ is bounded. More rigorously, if the limit is not infinity, there must exist a compact set $K \subset \R^N$ such that $\bv(t) \in K$ for all $t \in [0, T_+)$.

Because $K$ is compact and $f$ is locally Lipschitz, $f$ is uniformly Lipschitz on $K$ with some constant $L_K$, and $f$ is bounded on $K$, meaning there exists $M_K > 0$ such that $\|f(\bx)\| \leq M_K$ for all $\bx \in K$.
Since $\dot{\bv}(t) = f(\bv(t))$ and $\bv(t) \in K$, we have $\|\dot{\bv}(t)\| \leq M_K$. By the Mean Value Theorem, for any $t_1, t_2 \in [0, T_+)$:
\begin{equation}
\|\bv(t_1) - \bv(t_2)\| \leq M_K |t_1 - t_2|
\end{equation}

This shows that $\bv(t)$ is uniformly continuous on $[0, T_+)$. Because $\R^N$ is a complete metric space and $\bv(t)$ is uniformly continuous and bounded, the limit
\[
\bv_+ = \lim_{t \to T_+^-} \bv(t)
\]
must exist and is finite. We can now define a new initial value problem starting at time $T_+$ with the initial condition $\bv(T_+) = \bv_+$. By the local existence theorem (Picard--Lindelöf), there exists a solution $\tilde{\bv}(t)$ defined on some interval $[T_+, T_+ + \delta)$ for $\delta > 0$.
We can concatenate the original solution $\bv(t)$ and the new solution $\tilde{\bv}(t)$ to form a continuous, differentiable solution defined on $[0, T_+ + \delta)$. This strictly extends the domain of the solution past $T_+$, which directly contradicts the assumption that $(T_-, T_+)$ was the \textit{maximal} interval of existence.
Therefore, our assumption that the trajectory remains in a compact set must be false; it must escape to infinity.
\end{proof}

\begin{remark}[Practical Consequence]
The Escape Lemma provides a powerful practical tool: if you can prove that a trajectory is confined to a bounded region (for example, by finding a Lyapunov function or an invariant set), you immediately guarantee that $T_+ = \infty$, meaning the solution exists globally for all future time.
\end{remark}

\section{Stability and Local Dynamics}

\subsection{Lyapunov Stability}

When analyzing the long-term behavior of a system, we often study the behavior near an equilibrium point.

\begin{definition}[Equilibrium and Stability]
A point $\bv_\infty \in \R^N$ is an \textit{equilibrium} (or fixed point) if $f(\bv_\infty) = \mathbf{0}$. The constant function $\bv(t) = \bv_\infty$ is then a trivial solution.
The equilibrium $\bv_\infty$ is:
\begin{itemize}
\item \textit{Lyapunov stable}: if for every $\varepsilon > 0$, there exists a $\delta > 0$ such that if $\|\bv(0) - \bv_\infty\| < \delta$, then $\|\bv(t) - \bv_\infty\| < \varepsilon$ for all $t \geq 0$.
\item \textit{Asymptotically stable}: if it is Lyapunov stable, and there exists a $\delta_0 > 0$ such that if $\|\bv(0) - \bv_\infty\| < \delta_0$, then

\[
\lim_{t \to \infty} \bv(t) = \bv_\infty.
\]

\end{itemize}
\end{definition}

\begin{remark}[Geometric Meaning]
Lyapunov stability means that trajectories starting sufficiently close to the equilibrium remain arbitrarily close to it forever; they do not necessarily converge to it, but they never escape a small tube around it. Asymptotic stability adds the requirement that they actually converge to the equilibrium.
\end{remark}

To prove stability without explicitly solving the ODE, we use Lyapunov's Direct Method, which relies on energy-like scalar functions.

\begin{definition}[Lyapunov Function and Orbital Derivative]
Let $U$ be an open neighborhood of $\bv_\infty$. A continuously differentiable ($C^1$) function $V: U \to \R$ is a \textit{Lyapunov function} if:
\begin{enumerate}
\item \textit{Positive definiteness}: $V(\bv_\infty) = 0$ and $V(\bv) > 0$ for all $\bv \in U \setminus \{\bv_\infty\}$.
\item \textit{Orbital derivative}: the orbital derivative $\dot{V}(\bv) \leq 0$ for all $\bv \in U$.
\end{enumerate}
The \textit{orbital derivative} $\dot{V}(\bv)$ is the time derivative of $V$ evaluated along the trajectories of the system. By the multivariate chain rule:
\begin{equation}
\begin{aligned}
\dot{V}(\bv) &= \frac{\dd}{\dd t} V(\bv(t)) \bigg|_{t=0} \\
&= \nabla V(\bv) \cdot \dot{\bv} \\
&= \nabla V(\bv) \cdot f(\bv)
\end{aligned}
\end{equation}

\end{definition}

\begin{theorem}[Lyapunov's Direct Method]
If a Lyapunov function $V$ exists for $\bv_\infty$, then $\bv_\infty$ is Lyapunov stable. If, furthermore, $\dot{V}(\bv) < 0$ for all $\bv \in U \setminus \{\bv_\infty\}$ (strictly negative definite), then $\bv_\infty$ is asymptotically stable.
\end{theorem}

\begin{proof}[Rigorous Proof of Stability]
Let $\varepsilon > 0$ be given. We must find a $\delta > 0$ such that trajectories starting in the ball $B_\delta(\bv_\infty)$ never leave the ball $B_\varepsilon(\bv_\infty)$.
Consider the boundary of the $\varepsilon$-ball,
\[
S_\varepsilon = \{ \bv : \|\bv - \bv_\infty\| = \varepsilon \}.
\]

This set is closed and bounded, hence compact. Since $\bv_\infty \notin S_\varepsilon$ and $V$ is continuous and positive definite, $V$ attains a strictly positive minimum on $S_\varepsilon$. Let:
\begin{equation}
m = \min_{\bv \in S_\varepsilon} V(\bv) > 0
\end{equation}
Define the sublevel set
\[
U_m = \{ \bv \in U : V(\bv) < m \}.
\]

Because $V$ is continuous, $U_m$ is an open set. Since $V(\bv_\infty) = 0 < m$, $\bv_\infty \in U_m$. Therefore, there exists a $\delta > 0$ such that the ball $B_\delta(\bv_\infty) \subset U_m$.
Now, suppose $\bv(0) \in B_\delta(\bv_\infty)$. Then $V(\bv(0)) < m$.
Along the trajectory, the time derivative of $V$ is
\[
\begin{aligned}
(\dd/\dd t) V(\bv(t)) &= \dot{V}(\bv(t)) \\
&\leq 0.
\end{aligned}
\]
This means $V(\bv(t))$ is a monotonically non-increasing function of time. Thus, for all $t \geq 0$:
\begin{equation}
V(\bv(t)) \leq V(\bv(0)) < m
\end{equation}

This implies that $\bv(t) \in U_m$ for all $t \geq 0$. By the definition of $m$, any point on the boundary $S_\varepsilon$ has $V(\bv) \geq m$. Since $V(\bv(t)) < m$, the trajectory can never reach $S_\varepsilon$. Because the trajectory is continuous and starts inside $B_\varepsilon$, it can never cross $S_\varepsilon$ to exit the ball. Hence, $\|\bv(t) - \bv_\infty\| < \varepsilon$ for all $t \geq 0$, proving Lyapunov stability.
\end{proof}

\begin{proof}[Rigorous Proof of Asymptotic Stability]
We already know the equilibrium is stable, so trajectories starting sufficiently close remain close. We must now prove they converge to $\bv_\infty$.
Assume $\dot{V}(\bv) < 0$ for $\bv \neq \bv_\infty$. For a trajectory starting near $\bv_\infty$, $V(\bv(t))$ is monotonically decreasing and bounded below by $0$. By the Monotone Convergence Theorem for real sequences, the limit must exist:
\begin{equation}
\begin{aligned}
\lim_{t \to \infty} V(\bv(t)) &= c \\
&\geq 0
\end{aligned}
\end{equation}

We must show that $c = 0$. Suppose, for contradiction, that $c > 0$.
Because the trajectory is stable, it is confined to a compact set
\[
K = \{ \bv : c \leq V(\bv) \leq V(\bv(0)) \}.
\]

Note that $\bv_\infty \notin K$ because $V(\bv_\infty) = 0 < c$.
The orbital derivative $\dot{V}(\bv)$ is a continuous function on the compact set $K$. Since $\dot{V}(\bv) < 0$ everywhere on $K$, it must attain a strict maximum on $K$. Let this maximum be $-\eta$, where $\eta > 0$. Thus, for all $\bv \in K$:
\begin{equation}
\dot{V}(\bv) \leq -\eta < 0
\end{equation}
Evaluating this along the trajectory $\bv(t) \in K$:
\begin{equation}
\frac{\dd}{\dd t} V(\bv(t)) \leq -\eta
\end{equation}
Integrating both sides from $0$ to $t$:
\begin{equation}
V(\bv(t)) - V(\bv(0)) \leq -\eta t \implies V(\bv(t)) \leq V(\bv(0)) - \eta t
\end{equation}

As $t \to \infty$, the right-hand side approaches $-\infty$. This implies $V(\bv(t)) \to -\infty$, which strictly contradicts the fact that $V(\bv) \geq 0$ for all $\bv$.
Therefore, our assumption $c > 0$ must be false, and we conclude $c = 0$.
Since $V(\bv(t)) \to 0$ and $V$ is continuous and positive definite (meaning
\[
V(\bv) = 0 \iff \bv = \bv_\infty
\]
), it follows rigorously that
\[
\lim_{t \to \infty} \bv(t) = \bv_\infty.
\]

\end{proof}

\subsection{Hartman--Grobman and Center Manifolds}

When a Lyapunov function is difficult to construct, we often analyze the local behavior near an equilibrium by linearizing the system.

\begin{definition}[Jacobian and Hyperbolic Equilibrium]
The \textit{Jacobian matrix} of $f$ at $\bv_\infty$ is the $N \times N$ matrix $J$ of partial derivatives, with entries
\[
J_{ij} = (\partial f_i/\partial v_j)(\bv_\infty).
\]

The equilibrium $\bv_\infty$ is called \textit{hyperbolic} if every eigenvalue $\lambda$ of $J$ has a nonzero real part, i.e., $\Re(\lambda) \neq 0$.
\end{definition}

\begin{theorem}[Hartman--Grobman Theorem]
If $\bv_\infty$ is a hyperbolic equilibrium, then the nonlinear flow near $\bv_\infty$ is \textit{topologically conjugate} to the linear flow $\dot{\bepsilon} = J\bepsilon$.
\end{theorem}

\begin{remark}[Unpacking Topological Conjugacy]
What does ``topologically conjugate'' mean? It means there exists a homeomorphism (a continuous, bijective map with a continuous inverse) $h: U \to V$, where $U$ is a neighborhood of $\bv_\infty$ and $V$ is a neighborhood of the origin in $\R^N$, such that $h(\bv_\infty) = \mathbf{0}$, and $h$ maps trajectories of the nonlinear system to trajectories of the linear system, preserving the direction of time. Formally, if $\Phi_t(\bx)$ is the nonlinear flow and $e^{Jt}$ is the linear flow, then:
\begin{equation}
h(\Phi_t(\bx)) = e^{Jt} h(\bx)
\end{equation}

This implies that the phase portrait of the nonlinear system near a hyperbolic equilibrium is qualitatively identical to the phase portrait of its linearization. The topological structure (e.g., whether it is a node, saddle, or spiral) is completely determined by the eigenvalues of $J$.
\end{remark}

\begin{remark}[Why Hyperbolicity is Required]
If $J$ has an eigenvalue with zero real part ($\Re(\lambda) = 0$), the equilibrium is \textit{non-hyperbolic}. In this case, the linear terms are ``neutral'' (neither exponentially growing nor decaying), and the nonlinear higher-order terms in the Taylor expansion of $f$ can dominate the local dynamics. The nonlinear system could exhibit behaviors completely absent in the linearization, such as a center (closed orbits) or a weak focus. Thus, linearization is inconclusive for non-hyperbolic equilibria.
\end{remark}

To handle non-hyperbolic equilibria, we use the Center Manifold Theorem, which allows us to reduce the dimensionality of the system.

\begin{definition}[Spectral Subspaces]
Based on the eigenvalues of the Jacobian $J$, we decompose $\R^N$ into three invariant linear subspaces:
\begin{itemize}
\item \textit{Stable subspace $E^s$}: the span of all generalized eigenvectors corresponding to eigenvalues with $\Re(\lambda) < 0$.
\item \textit{Unstable subspace $E^u$}: the span of all generalized eigenvectors corresponding to eigenvalues with $\Re(\lambda) > 0$.
\item \textit{Center subspace $E^c$}: the span of all generalized eigenvectors corresponding to eigenvalues with $\Re(\lambda) = 0$.
\end{itemize}
By linear algebra, $\R^N = E^s \oplus E^u \oplus E^c$.
\end{definition}

\begin{theorem}[Center Manifold Theorem]
There exists a locally invariant manifold $W^c$ (the center manifold) passing through $\bv_\infty$ such that:
\begin{enumerate}
\item \textit{Tangency}: $W^c$ is tangent to the center subspace $E^c$ at $\bv_\infty$.
\item \textit{Local invariance}: $W^c$ is locally invariant: if a trajectory starts in $W^c$, it remains in $W^c$ for all time (as long as it stays in the neighborhood).
\item \textit{Dimension}: the dimension of $W^c$ is equal to the dimension of $E^c$.
\end{enumerate}
Furthermore, there exist locally invariant stable and unstable manifolds $W^s$ and $W^u$ tangent to $E^s$ and $E^u$, respectively.
\end{theorem}

\begin{remark}[Why is this Useful?]
The dynamics on the stable manifold $W^s$ are trivial: trajectories exponentially decay to $\bv_\infty$. The dynamics on the unstable manifold $W^u$ are also well-understood: trajectories exponentially repel from $\bv_\infty$.
The complex, interesting dynamics (bifurcations, stability transitions, limit cycles) occur entirely on the center manifold $W^c$. Because
\[
\dim(W^c) = \dim(E^c) < N
\]

 (assuming $E^s$ or $E^u$ are non-trivial), the Center Manifold Theorem allows us to reduce the study of the full $N$-dimensional nonlinear system to a lower-dimensional system restricted to $W^c$. On this reduced space, the linear terms are zero, and the dynamics are governed purely by the nonlinear terms, which can be analyzed using techniques like Lyapunov functions or normal forms.
\end{remark}

Figure~\ref{fig:ch15-phase-portraits} compares the local phase portraits of four types of planar equilibrium.

\begin{figure}[!htbp]
\centering
\begin{tikzpicture}
\begin{axis}[width=6.1cm,height=6.1cm,axis lines=middle,xmin=-2,xmax=2,ymin=-2,ymax=2,
  xtick=\empty,ytick=\empty,]
  \foreach \sx/\sy in {1/1,-1/1,1/-1,-1/-1}{
    \foreach \c in {0.4,0.8,1.4,2.2}{
      \addplot[figblue,thick,domain=0:2.2,samples=50,forget plot]
        ({\sx*\c*exp(-0.6*x)},{\sy*\c*0.4*exp(-0.3*x)});
    }
    \addplot[figblue,thick,-{Stealth[length=2.6mm]},forget plot] coordinates {
      ({\sx*1.4*exp(-0.6*1.55)},{\sy*1.4*0.4*exp(-0.3*1.55)})
      ({\sx*1.4*exp(-0.6*1.85)},{\sy*1.4*0.4*exp(-0.3*1.85)})
    };
  }
\end{axis}
\begin{axis}[at={(6.6cm,0)},width=6.1cm,height=6.1cm,axis lines=middle,xmin=-2,xmax=2,ymin=-2,ymax=2,
  xtick=\empty,ytick=\empty,]
  \foreach \c in {0.3,0.7,1.2,2}{
    \addplot[figred,thick,domain=-1.5:1.5,samples=60,forget plot] ({x},{\c/(x+0.001)});
    \addplot[figred,thick,domain=-1.5:1.5,samples=60,forget plot] ({x},{-\c/(x+0.001)});
  }
  \draw[vecB] (-1.9,-1.9) -- (-0.35,-0.35);
  \draw[vecB] (1.9,1.9) -- (0.35,0.35);
  \draw[vecB] (0.35,-0.35) -- (1.9,-1.9);
  \draw[vecB] (-0.35,0.35) -- (-1.9,1.9);
\end{axis}
\begin{axis}[at={(0,-6.6cm)},width=6.1cm,height=6.1cm,axis lines=middle,xmin=-2,xmax=2,ymin=-2,ymax=2,
  xtick=\empty,ytick=\empty,]
  \addplot[figgreen,thick,domain=0:12,samples=200,forget plot,variable=t]
    ({1.9*exp(-0.15*t)*cos(deg(t))},{1.9*exp(-0.15*t)*sin(deg(t))});
  \addplot[figgreen,thick,-{Stealth[length=2.6mm]},forget plot] coordinates {
    ({1.9*exp(-0.15*7.4)*cos(deg(7.4))},{1.9*exp(-0.15*7.4)*sin(deg(7.4))})
    ({1.9*exp(-0.15*7.8)*cos(deg(7.8))},{1.9*exp(-0.15*7.8)*sin(deg(7.8))})
  };
\end{axis}
\begin{axis}[at={(6.6cm,-6.6cm)},width=6.1cm,height=6.1cm,axis lines=middle,xmin=-2,xmax=2,ymin=-2,ymax=2,
  xtick=\empty,ytick=\empty,]
  \foreach \r in {0.5,0.9,1.3,1.7}{
    \addplot[black!55,thick,domain=0:360,samples=80,forget plot] ({\r*cos(x)},{\r*sin(x)});
  }
\end{axis}
\end{tikzpicture}
\caption[Phase portraits: node, saddle, focus, center]{Local phase portraits around the origin. Top left: stable node with two negative real eigenvalues. Top right: saddle with one negative and one positive real eigenvalue. Bottom left: stable focus with complex eigenvalues having negative real part. Bottom right: non-hyperbolic center with purely imaginary eigenvalues. Arrows indicate the direction of evolution where shown.}
\label{fig:ch15-phase-portraits}
\end{figure}

\section{Partial Differential Equations: Solution Methods}
\label{sec:pde-methods}

The ordinary differential equations studied so far describe finitely many state variables evolving in time. Spatially distributed systems require an unknown field $u(\bx,t)$ and partial derivatives in both space and time. The integral theorems, spectral methods, and distribution theory developed in earlier chapters now provide tools for solving these equations.

\subsection{Fundamental PDEs and Separation of Variables}

Three fundamental linear second-order partial differential equations are
\begin{align}
\text{Heat (diffusion) equation:}\qquad &\frac{\partial u}{\partial t}=D\,\Delta u,\\
\text{Wave equation:}\qquad &\frac{\partial^2u}{\partial t^2}=c^2\,\Delta u,\\
\text{Laplace (steady-state) equation:}\qquad &\Delta u=0,
\end{align}
Here the Laplacian is
\[
\begin{aligned}
\Delta u &= \nabla\cdot(\nabla u) \\
&= \sum_i\partial^2u/\partial x_i^2,
\end{aligned}
\]

The constant $D>0$ is the diffusion coefficient, and $c>0$ is the propagation speed. In mathematical biology, the diffusion equation models the spatial spreading of a cellular or molecular concentration. Initial and boundary conditions are needed to select a solution appropriate to the physical problem.


The basic idea is to seek a factored solution
\[
u(\bx,t)=X(\bx)T(t).
\]

Substitution into a suitable linear PDE separates the spatial and temporal dependence, linking them by a separation constant. In one spatial dimension, this produces two ordinary differential equations; in higher dimensions, the spatial factor generally satisfies an eigenvalue PDE.

\begin{example}[The Heat Equation with Homogeneous Dirichlet Conditions]
\label{ex:heat-separation}
Consider $u_t=D\,u_{xx}$ for $0<x<L$ and $t>0$, with
\[
u(0,t)=u(L,t)=0,\qquad u(x,0)=f(x).
\]

The boundary conditions are called \emph{homogeneous Dirichlet conditions} because the function vanishes at the boundary.

\textit{Step 1: separate the variables.} Substituting

\[
u(x,t)=X(x)T(t)
\]
gives $XT'=DX''T$. Wherever the factors are nonzero,
\begin{equation}
\frac{T'(t)}{D\,T(t)}=\frac{X''(x)}{X(x)}.
\end{equation}

The left-hand side depends only on $t$ and the right-hand side only on $x$. They must therefore equal a common constant, denoted by $-\lambda$:
\begin{equation}
T'=-\lambda D\,T,\qquad X''=-\lambda X.
\end{equation}

These equations are then imposed throughout the interval, including any zeros of $X$; the quotient is only a device for identifying the separated equations.

\textit{Step 2: solve the spatial eigenvalue problem.} A nonzero temporal factor requires

\[
\begin{aligned}
X(0) &= X(L) \\
&= 0.
\end{aligned}
\]
For $\lambda>0$,
\[
X(x)=A\cos(\sqrt\lambda x)+B\sin(\sqrt\lambda x).
\]

The first boundary condition gives $A=0$, and the second requires $\sin(\sqrt{\lambda} L)=0$ for a nontrivial solution. Hence
\begin{equation}
\lambda_n=\left(\frac{n\pi}{L}\right)^2,
\qquad X_n(x)=\sin\left(\frac{n\pi x}{L}\right),\qquad n=1,2,\dots.
\end{equation}
No nontrivial solution has $\lambda\le0$: integration by parts in $-X''=\lambda X$ gives
\[
\int_0^L|X'|^2\,\dd x=\lambda\int_0^L|X|^2\,\dd x,
\]
and the zero endpoint values rule out a nonzero constant function.

\textit{Step 3: solve the temporal equation and superpose.} For each eigenvalue,

\[
T_n(t)=e^{-\lambda_nDt}
\]
up to a multiplicative constant. Linearity permits the expansion
\begin{equation}
u(x,t)=\sum_{n=1}^{\infty}b_n\sin\left(\frac{n\pi x}{L}\right)e^{-\lambda_nDt}.
\end{equation}

Matching the initial data amounts to expanding $f$ in a Fourier sine series. Orthogonality on $(0,L)$ gives
\begin{equation}
b_n=\frac{2}{L}\int_0^L f(x)\sin\left(\frac{n\pi x}{L}\right)\,\dd x.
\end{equation}

For $f\in L^2(0,L)$, the initial condition is attained in the $L^2$ sense as $t\downarrow0$, and the exponential factors smooth the series for $t>0$. With suitable smoothness and endpoint compatibility, it gives the classical solution up to the initial boundary. The functions $X_n$ are eigenfunctions of the self-adjoint Dirichlet operator $-\dd^2/\dd x^2$, just as orthogonal eigenvectors arise for symmetric matrices. Thus separation of variables connects ODE methods, Fourier expansions, and operator theory.
\end{example}

Figure~\ref{fig:atlas-heat-mode-decay} illustrates the different decay rates of separated heat-equation modes.

\begin{figure}[!htbp]
\centering
\begin{tikzpicture}[>=Stealth]
\begin{axis}[width=11cm,height=5.8cm,axis lines=middle,ticklabel style={font=\small},label style={font=\small},legend style={font=\scriptsize,draw=black!20,fill=white},xlabel={$x$},ylabel={$u(x,t)$},xmin=0,xmax=1.05,ymin=0,ymax=1.3,xtick={0,0.5,1},legend columns=3,legend pos=north east]
\addplot[figblue,very thick,domain=0:1,samples=150] {sin(deg(pi*x))+0.35*sin(deg(3*pi*x))};
\addlegendentry{$t=0$}
\addplot[figred,thick,dashed,domain=0:1,samples=150] {exp(-pi^2*0.03)*sin(deg(pi*x))+0.35*exp(-9*pi^2*0.03)*sin(deg(3*pi*x))};
\addlegendentry{$t=0.03$}
\addplot[figgreen,thick,dashdotted,domain=0:1,samples=150] {exp(-pi^2*0.1)*sin(deg(pi*x))+0.35*exp(-9*pi^2*0.1)*sin(deg(3*pi*x))};
\addlegendentry{$t=0.1$}
\end{axis}
\end{tikzpicture}
\caption{Three profiles solve $u_t=u_{xx}$ on $[0,1]$ with zero boundary values. The initial profile combines the first sine mode with $0.35$ times the third. The third mode decays nine times faster in the exponent, so small-scale structure disappears before the main profile relaxes to zero.}
\label{fig:atlas-heat-mode-decay}
\end{figure}

\subsection{Dirichlet and Neumann Boundary-Value Problems}
\label{sec:pde-boundary-problems}

\begin{definition}[Dirichlet and Neumann Problems]
Let $\Omega\subset\R^N$ be a bounded connected domain with smooth boundary. A \emph{Dirichlet problem} for the Laplace equation prescribes the boundary values:
\begin{equation}
\Delta u=0\quad\text{in }\Omega,\qquad u=g\quad\text{on }\partial\Omega.
\end{equation}
A \emph{Neumann problem} instead prescribes the outward normal derivative:
\begin{equation}
\Delta u=0\quad\text{in }\Omega,\qquad
\frac{\partial u}{\partial\mathbf{n}}=h\quad\text{on }\partial\Omega,
\end{equation}
where $\partial u/\partial\mathbf{n}:=\nabla u\cdot\mathbf{n}$. This specifies a quantity proportional to the boundary flux rather than the value of $u$ itself; for a diffusive current $\mathbf{J}=-D\nabla u$, the outward flux is $-D h$.
\end{definition}

\begin{theorem}[Uniqueness for the Dirichlet Problem]
If a classical solution of the Dirichlet problem exists with sufficient regularity for integration by parts, it is unique.
\end{theorem}

\begin{proof}[The Energy Method]
Let $u_1,u_2$ be two solutions and set $w=u_1-u_2$. Then $\Delta w=0$ in $\Omega$ and $w=0$ on $\partial\Omega$. The product rule gives
\begin{equation}
\begin{aligned}
\nabla\cdot(w\nabla w) &= |\nabla w|^2+w\Delta w \\
&= |\nabla w|^2.
\end{aligned}
\end{equation}
Integrating and applying the $N$-dimensional version of the divergence theorem yields
\begin{equation}
\int_\Omega|\nabla w|^2\,\dd\bx
=\int_{\partial\Omega}w\frac{\partial w}{\partial\mathbf{n}}\,\dd S=0.
\end{equation}

The continuous nonnegative integrand must vanish identically. Thus $\nabla w=0$, and connectedness implies that $w$ is constant. Its zero boundary value forces that constant to be zero, so $u_1=u_2$.
\end{proof}

\begin{remark}[Neumann Uniqueness and Compatibility]
For two solutions of the same Neumann problem, the difference satisfies $\partial w/\partial\mathbf{n}=0$, which again makes the energy boundary term vanish. Therefore $w$ is constant, but the constant need not be zero: the solution is unique only up to an additive constant. Existence also requires the compatibility condition
\begin{equation}
\int_{\partial\Omega}h\,\dd S=\int_\Omega\Delta u\,\dd\bx=0.
\end{equation}
Prescribing, for example, the mean of $u$ fixes the additive constant.
\end{remark}

\subsection{Green's Functions}
\label{sec:green-functions}

We now consider Poisson's equation $\Delta u=-\rho$ in $\Omega$, with homogeneous Dirichlet data. The idea parallels the use of the Dirac delta in distribution theory: if the response to a point source at $\by$ is known, the response to a distributed source $\rho$ can be reconstructed by integrating these point-source responses, weighted by $\rho(\by)$.

\begin{definition}[Dirichlet Green's Function]
A \emph{Green's function} for $-\Delta$ on $\Omega$ with homogeneous Dirichlet conditions satisfies, for each fixed $\by\in\Omega$,
\begin{equation}
-\Delta_{\bx}G(\bx,\by)=\delta(\bx-\by)\quad\text{in }\Omega,
\qquad G(\bx,\by)=0\quad\text{on }\partial\Omega.
\end{equation}

The differential equation is interpreted in the distributional sense. The multidimensional Dirac delta acts on a test function by $\delta_{\by}[\varphi]=\varphi(\by)$.
\end{definition}

\begin{theorem}[Green's Function Representation]
For a sufficiently regular domain and source for which the Dirichlet Green's function representation is valid, the solution of $\Delta u=-\rho$ in $\Omega$, $u=0$ on $\partial\Omega$, is
\begin{equation}
u(\bx)=\int_\Omega G(\bx,\by)\rho(\by)\,\dd\by.
\end{equation}

\end{theorem}

\begin{proof}[Distributional Verification]
Formally applying the Laplacian gives
\begin{align}
\Delta_{\bx}u(\bx)
&=\int_\Omega\Delta_{\bx}G(\bx,\by)\rho(\by)\,\dd\by\notag\\
&=-\int_\Omega\delta(\bx-\by)\rho(\by)\,\dd\by=-\rho(\bx).
\end{align}

Because $G$ is singular at $\bx=\by$, differentiation here is understood distributionally rather than as an unqualified exchange of classical derivatives and integration. More explicitly, for a compactly supported smooth test function $\varphi$, the defining identity for $G$ and Fubini's theorem give
\begin{align}
\int_\Omega u(\bx)(-\Delta\varphi(\bx))\,\dd\bx
&=\int_\Omega\rho(\by)\left[\int_\Omega G(\bx,\by)(-\Delta\varphi(\bx))\,\dd\bx\right]\dd\by\notag\\
&=\int_\Omega\rho(\by)\varphi(\by)\,\dd\by.
\end{align}

Under the assumed regularity, the boundary trace is zero because the Green's function has zero Dirichlet trace. When the represented solution is classical, the preceding uniqueness theorem identifies it as the only solution.
\end{proof}

\begin{example}[The Upper Half-Plane and the Method of Images]
Consider the upper half-plane
\[
\Omega=\{\bx=(x_1,x_2):x_2>0\}\subset\R^2
\]

Let $\by=(y_1,y_2)$ with $y_2>0$. Reflect the source across the boundary to obtain $\by^*=(y_1,-y_2)$. Then
\begin{equation}
G(\bx,\by)=\frac{1}{2\pi}\log\frac{|\bx-\by^*|}{|\bx-\by|}
\end{equation}
is a Dirichlet Green's function for the upper half-plane. Indeed, the two-dimensional fundamental solution satisfies
\[
-\Delta_{\bx}\left(-\frac{1}{2\pi}\log|\bx-\by|\right)=\delta(\bx-\by).
\]

The added term $\log|\bx-\by^*|/(2\pi)$ is harmonic throughout $\Omega$ because its singularity lies outside the domain. On $x_2=0$, the distances to $\by$ and $\by^*$ coincide, so the logarithmic ratio vanishes. The normalization of the fundamental solution follows from the divergence theorem on a small circle around the source, just as $1/(4\pi|\bx-\by|)$ is the fundamental solution of $-\Delta$ in three dimensions.

Adding a fictitious exterior source to enforce the boundary condition is the \emph{method of images}. It is particularly useful on domains with suitable symmetries, such as half-spaces and balls, although the image location and strength depend on the geometry. On an unbounded domain, appropriate behavior at infinity must also be specified to obtain uniqueness.
\end{example}